% chapter_04_MPI.tex -- chapter 04 of the JitJet book: multiparton interactions and beam remnants.
% Script: test/04_MPI/run_04_MPI.sh; every number is from test/04_MPI/out/ (ledger_04_MPI.txt, genPartAnalyzer_04_MPI.log,
% step1_GEN_04_MPI.log) or from FWLite on test/04_MPI/out/step1_GEN_04_MPI.root (companion chapter_04_MPI.md, section 4).
\chapter{04 -- MPI: the other partons of the two protons scatter too, and the proton remnants are built}
\label{ch:04_MPI}
\providecommand{\ptz}{\ensuremath{p_{\mathrm{T}0}}}  % the MPI regulator

\section{What this chapter is about}

A proton is not one parton. When the $\bar u$ of one proton and the $u$ of the other produce our \bbbar{} pair
(chapter~\ref{ch:01_ME}), the two protons pass through each other as a whole, and nothing forbids a second,
third, \dots{} pair of partons from scattering in the same collision. Pythia8 calls these extra scatterings
\emph{multiparton interactions} (MPI); their products, together with the radiation they emit, form the
\emph{underlying event} that every jet measurement has to live with. This chapter switches them on
(\code{PartonLevel:MPI = on}) and lets Pythia8 finish its \emph{parton level}: the MPI, the \emph{primordial
$k_\mathrm{T}$} given to every incoming parton, and the \emph{beam remnants}, the partons of the two protons that
took part in nothing and still carry most of the $13.6\TeV$. Hadronisation and decays stay off
(\code{HadronLevel:Hadronize = off}, \code{HadronLevel:Decay = off}), so the output is still a record of quarks and
gluons.

The event grows from 11 final-state partons (chapter~\ref{ch:02_PS}) to 169: Pythia8 generated 31~MPI besides
the hard process. This chapter counts them, looks at how hard they are, where they go, how much of them lands near
the two $b$ jets, and what the regulator $\ptz$ of the CP5 tune has to do with it. From here on the event is
final at parton level: chapters~\ref{ch:05_Had} and~\ref{ch:06_Decay} build on exactly this record.

\begin{plainwords}[Multiparton interactions in one sentence]
Every $2\to2$ QCD scattering has a cross section that grows like $1/\pt^4$ towards small \pt, so at some
\pt{} the parton--parton cross section exceeds the proton--proton one: the only consistent reading is that
several parton pairs scatter in the same $pp$ collision. Pythia8 generates them one after the other in
decreasing \pt, with the hard process ($\ptHat = 113\GeV$) first, and damps the divergence with a regulator,
$\alpha_s(\pt^2)/\pt^4 \to \alpha_s(\pt^2+\ptz^2)/(\pt^2+\ptz^2)^2$. The number of MPI per collision is
what the tune (here CP5) is fitted to reproduce; a minimum-bias collision (chapter~\ref{ch:08_PU}) is the same
mechanism without any hard process on top.
\end{plainwords}

\section{From the previous chapter}

Chapter~\ref{ch:03_MnM} changed nothing (a leading-order $2\to2$ process has nothing to match or merge), so the
input of this stage is the configuration of chapter~\ref{ch:02_PS} with one more switch,
\code{PartonLevel:MPI = on}; same fragment, same seeds, same conditions. What is kept is the hard process:
the four entries of status 21--23 are identical in every generator chapter (incoming $\bar u$ with $x=0.0022$
and $u$ with $x=0.3901$, $\sqrt{\hat s}=396.1\GeV$, $b$ and $\bar b$ at $\pt=113.2\GeV$, $\eta=-1.44$ and
$-3.75$).\footnote{\code{ledger\_04\_MPI.txt}, identical lines in \code{ledger\_02\_PS.txt}; the Pythia8 listing of the
hard process is in \code{step1\_GEN\_04\_MPI.log}.}

What changes is everything after the hard process, including the parton shower of chapter~\ref{ch:02_PS}, and
figure~\ref{fig:04_showers} shows it at once: the display after the shower does \emph{not} match the one of
chapter~\ref{ch:02_PS}, and it cannot. Pythia8 does not run the showers and then the MPI: it runs ISR, FSR and MPI
in one sequence of decreasing \pt{} (the \emph{interleaved evolution}), each candidate emission or interaction
competing for the next step. With MPI on, the random numbers are consumed in a different order from the first MPI
trial on, and every later emission is a different draw; Pythia8 has no switch that adds MPI on top of a shower
that already exists. Table~\ref{tab:04_shower} gives the numbers: the shower of chapter~\ref{ch:02_PS} had no ISR
and five FSR emissions, the $b$ lost more than half of its \pt{} to its radiation ($113.2\to46.9\GeV$) and the
$\bar b$ did not radiate at all; the event of this chapter has 28~ISR and 64~FSR emissions, but only 2 and 10 of
them in the shower of the hard process (the others are the showers of the 31~MPI, drawn with them at stage~3), and
the $b$ ends at $91.6\GeV$ and the $\bar b$ at $118.6\GeV$. Neither is ``more right''; they are two draws of the
same probability distribution. The shower of this chapter is the one that stays: chapters~\ref{ch:05_Had}
and~\ref{ch:06_Decay} are identical to it at parton level, and the 14 partons of figure~\ref{fig:04_showers}(b)
are found again, moved only by the primordial $k_\mathrm{T}$, in the display at the end of this chapter
(figure~\ref{fig:04_display}).

The dashed circles are the cones $\Delta R<0.4$ around the $b$ and the $\bar b$, what a jet of radius 0.4 would
collect, drawn as in chapter~\ref{ch:02_PS}. There the $b$ cone held the $b$ and its five FSR partons, exactly
the $113.2\GeV$ of the hard-process $b$. Here it holds $120.3\GeV$ in six partons, more than the $b$ had: the
$b$ of $91.6\GeV$ and five partons of 15.3, 5.3, 3.5, 2.9 and $1.7\GeV$ that all descend from the first ISR
emission of the hard process, a gluon of $8.6\GeV$ emitted at $\eta=-1.93$, $\phi=-1.55$ ($\Delta R=0.64$ from
the $b$), whose own radiation (drawn as FSR) ends inside the cone of the $b$. A cone does not know whether a
parton was radiated before or after the collision. The second ISR gluon, $12.0\GeV$ at $\eta=1.6$ (the orange
triangle), lands far from both cones, and the $\bar b$ cone holds the $\bar b$ alone, $118.6\GeV$.%
\footnote{\code{genPartAnalyzer\_04\_MPI.log}, ``cone'' lines of stage 2; the ancestry of the five partons
from FWLite on \code{step1\_GEN\_04\_MPI.root} (companion, section 4.6).}

\begin{figure}[tb]
\centering
\begin{minipage}[t]{0.495\textwidth}\centering
{\small\textbf{(a)} chapter~\ref{ch:02_PS}: MPI off}\\[2pt]
\includegraphics[width=\textwidth,trim=0 0 170 0,clip]{02_PS/out/genPartAnalyzer_02_PS_page-2.png}
\end{minipage}\hfill
\begin{minipage}[t]{0.495\textwidth}\centering
{\small\textbf{(b)} this chapter: MPI on, same hard process}\\[2pt]
\includegraphics[width=\textwidth,trim=0 0 170 0,clip]{04_MPI/out/genPartAnalyzer_04_MPI_page-2.png}
\end{minipage}
\caption{The parton shower of the same hard process twice, at the end of stage 2 (page 2 of
\code{genPartAnalyzer\_02\_PS.pdf} and of \code{genPartAnalyzer\_04\_MPI.pdf} without the legend panel; (a) in
full is figure~\ref{fig:02_PS_display}). Page 1, the hard process, is the same picture in both. (a) MPI off:
the $b$ at $46.9\GeV$ with five FSR partons (green triangles), no ISR, the $\bar b$ untouched at $113.2\GeV$,
$b+\bar b$ $66.4\GeV$. (b) MPI on: the $b$ at $91.6\GeV$, the $\bar b$ at $118.6\GeV$, one ISR parton (orange,
$12.0\GeV$) and 11 FSR partons, $b+\bar b$ $27.3\GeV$; only the hard process and its own showers are drawn, the
31 MPI and their showers come at stage 3 (figure~\ref{fig:04_display}). Dashed: the cones $\Delta R<0.4$, 113.2
and $113.2\GeV$ in (a), 120.3 and $118.6\GeV$ in (b); the insets magnify the $b$ cone. Same seeds: two draws of
the interleaved evolution.}
\label{fig:04_showers}
\end{figure}

\begin{table}[htb]
\centering\small
\caption{The parton shower of this event in chapter~\ref{ch:02_PS} (MPI off) and in this chapter (MPI on), same
hard process, same seeds. The $b$ and $\bar b$ are their last shower copies (status 52; status 23 for the
$\bar b$ of chapter~\ref{ch:02_PS}, which did not radiate). Emission counts from the Pythia8 ``Info Listing''
of \code{step1\_GEN\_02\_PS.log} and \code{step1\_GEN\_04\_MPI.log} (its ``Number of MPI'' includes the hard
process; they count the emissions of the whole event), the rest from \code{genPartAnalyzer\_02\_PS.log} and
\code{genPartAnalyzer\_04\_MPI.log}; the split by scattering (each emission followed to the scattering it belongs
to) from FWLite on \code{step1\_GEN\_04\_MPI.root}. The analyzer draws at stage 2 only the hard process and its
own showers (figure~\ref{fig:04_showers}); the emissions of an MPI shower, whose chain of incoming partons leads to
an MPI, come with stage 3. The cones are the analyzer's (option \code{coneR = 0.4}): the partons drawn at stage 2
within $\Delta R<0.4$ of the quark, their number and vector \pt.}
\label{tab:04_shower}
\begin{tabular}{lrrr}
\toprule
after the parton shower (stage 2) & chapter 02 & chapter 04 & difference \\
\midrule
ISR emissions & 0 & 28 & $+28$ \\
\quad of which in the shower of the hard process & 0 & 2 & $+2$ \\
FSR emissions & 5 & 64 & $+59$ \\
\quad of which in the shower of the hard process & 5 & 10 & $+5$ \\
MPI (Pythia8 count, hard process included) & 1 & 32 & $+31$ \\
partons drawn at stage 2 (the hard process and its showers) & 7 & 14 & $+7$ \\
$b$: \pt{} [GeV] & 46.88 & 91.59 & $+44.71$ \\
$\bar b$: \pt{} [GeV] & 113.22 & 118.57 & $+5.35$ \\
$|\vec\pt(b)+\vec\pt(\bar b)|$ [GeV] (recoil against the radiation) & 66.40 & 27.29 & $-39.12$ \\
$b$ cone $\Delta R<0.4$: $n$ / \pt{} [GeV] & 6 / 113.22 & 6 / 120.27 & 0 / $+7.05$ \\
$\bar b$ cone $\Delta R<0.4$: $n$ / \pt{} [GeV] & 1 / 113.22 & 1 / 118.57 & 0 / $+5.35$ \\
\bottomrule
\end{tabular}
\end{table}

\begin{plainwords}[Interleaved evolution]
Pythia8 orders ISR, FSR and MPI in one falling \pt{} scale, starting at the hard scale ($113.3\GeV$ here, the
``Max \pt{} scale'' of the log) and stopping at the shower cut-off and at the MPI regulator. At each step the
three candidates propose a \pt, the highest wins, and the event is updated before the next step. The MPI
switch changes the shower not because the physics of ISR or FSR changed but because the sequence of random draws
did; only the hard process, chosen before the evolution starts, is untouched.
\end{plainwords}

\section{The numbers of the event}

\paragraph{How many interactions, and how hard.}
The MPI partons carry status 31 (incoming) and 33 (outgoing) in the record; the 124 entries of the group
31--39 are 62 of each, two incoming and two outgoing partons per interaction, so this collision has 31~MPI on top
of the hard process, the ``Number of MPI = 32'' of the Pythia8 log.\footnote{\code{ledger\_04\_MPI.txt} (status
groups), \code{step1\_GEN\_04\_MPI.log} (Info Listing). The status-31 count is the method of
\code{test/common/minBiasGenStats.py} for the pileup library of chapter~\ref{ch:08_PU}.} The same log explains
why so many: the impact parameter of the collision is small, $b=0.31$ in Pythia8's units, which gives an
``enhancement factor'' of 3.709 for the MPI rate; a hard process selects central collisions, where the two
protons overlap most. For comparison, the 250 minimum-bias collisions of chapter~\ref{ch:08_PU} have 7.0
interactions on average (median 3, maximum 39).\footnote{\code{test/08\_PU/out/genstats\_08\_PU.txt}.} The
interactions are soft: figure~\ref{fig:04_mpipt} gives the \pt{} of the outgoing pair of each of the 31, from
$6.65\GeV$ ($s\,g\to s\,g$) down to $0.24\GeV$; 18 of the 31 are below $1\GeV$ and only five above $4\GeV$.
Their 62 outgoing partons are 39 gluons, 11 $u$, 7 $d$, 4 $s$ and one $c$ (quark or antiquark), from incoming
partons with momentum fractions $x$ between $2.9\times10^{-7}$ and $0.24$ (median $2.9\times10^{-3}$): the
sea of the two protons at small $x$, where gluons dominate.\footnote{FWLite on \code{step1\_GEN\_04\_MPI.root};
the list is in section 4 of the companion.}

\begin{figure}[htb]
\centering
\begin{tikzpicture}[x=1.6cm, y=0.22cm]
\draw[->] (0,0) -- (5.4,0) node[right, font=\small] {$\pt$ of the pair [GeV]};
\draw[->] (0,0) -- (0,14) node[above, font=\small] {interactions};
\histbar{0}{7}\histbar{1}{11}\histbar{2}{3}\histbar{3}{5}\histbar{4}{5}
\foreach \x/\l in {0/$<0.5$, 1/0.5--1, 2/1--2, 3/2--4, 4/$>4$} \node[below, font=\small] at (\x+0.5,0) {\l};
\foreach \y in {5,10} \draw (0.1,\y) -- (-0.1,\y) node[left, font=\scriptsize] {\y};
\end{tikzpicture}
\caption{The 31 multiparton interactions of the event by the \pt{} of their outgoing parton pair (FWLite on
\code{step1\_GEN\_04\_MPI.root}, status-33 partons grouped by their mothers). The regulator of the CP5 tune at
$13.6\TeV$ is $\ptz=1.44\GeV$; the generation stops at \code{MultipartonInteractions:pTmin} $=0.2\GeV$
(Pythia8 default).}
\label{fig:04_mpipt}
\end{figure}

\paragraph{The regulator $\ptz$ and its energy dependence.}
The MPI cross section is regulated by $\ptz$, which Pythia8 scales with the collision energy,
\[
\ptz(\sqrt s)=\ptz^{\mathrm{Ref}}\left(\frac{\sqrt s}{E_{\mathrm{Ref}}}\right)^{\epsilon},
\qquad
\ptz^{\mathrm{Ref}}=1.41\GeV,\quad E_{\mathrm{Ref}}=7\TeV,\quad \epsilon=0.03344 ,
\]
with \code{MultipartonInteractions:pT0Ref} and \code{ecmPow} from the CP5 block
(\path{Configuration/Generator/python/MCTunes2017/PythiaCP5Settings_cfi.py}) and \code{ecmRef} at its Pythia8
default. This gives $\ptz=1.4395\GeV$ at $13\TeV$, $1.4417\GeV$ at $13.6\TeV$ (the ``pT0 = 1.44'' of the
initialisation in the log) and $1.4431\GeV$ at $14\TeV$: a $0.1\%$ step from Run~3 to the HL-LHC. The same
initialisation prints what the regulator does: with $\ptz=1.44\GeV$ the integrated parton--parton cross section
above the cut-off is $\sigma_{\mathrm{int}}=500.67\mb$ against a non-diffractive $pp$ cross section
$\sigma_{\mathrm{ND}}=55.42\mb$; their ratio, 9.0, is the mean number of interactions per non-diffractive
collision before the impact-parameter weighting. The other CP5 MPI parameters are the matter profile
(\code{bProfile=2}, a double Gaussian with \code{coreRadius=0.7634} and \code{coreFraction=0.63}) and
$\alpha_s(M_Z)=0.118$ at NLO for the MPI, the showers and the hard process.

\begin{plainwords}[What a tune is]
Pythia8 cannot compute $\ptz$, the proton's matter profile or how often colour lines
reconnect; these are parameters fitted to data. CP5 (\code{Tune:pp 14}, Monash 2013, with the CMS changes of the
block above) was fitted by CMS in 2017 (CMS-GEN-17-001) to underlying-event and minimum-bias data at 7 and
$13\TeV$ with NNPDF3.1 at NNLO. Everything in this chapter except the hard process depends on them; the ``Up''/``Down'' variants of the block
(\code{pT0Ref} 1.407 and 1.46) are the tune uncertainty.
\end{plainwords}

\paragraph{Primordial $k_\mathrm{T}$.}
The partons inside a proton are not exactly collinear with it. Once the interleaved evolution is finished,
Pythia8 gives every incoming parton of every interaction a small transverse kick and takes the recoil from the
remnants: the 64 entries of status 61 are these 64 incoming partons, two per interaction, with a mean
$k_\mathrm{T}$ of $1.09\GeV$ (median $0.77\GeV$, largest $3.89\GeV$).\footnote{FWLite on
\code{step1\_GEN\_04\_MPI.root}, status-61 entries.} The outgoing partons of each interaction are then copied
with the kick included: the $b$ goes from $\pt=91.59$ to $91.74\GeV$, the $\bar b$ from 118.57 to
$116.36\GeV$, and the net $|\vec\pt(b)+\vec\pt(\bar b)|$ from 27.29 to $25.08\GeV$. In this file these final
copies have no daughters, so they are the status-1 partons of the record, the ones page~3 of the analyzer draws;
in chapter~\ref{ch:05_Had}, where they fragment, the same copies appear as status 62.

\begin{figure}[!tb]
\centering
\includegraphics[width=\textwidth]{04_MPI/out/genPartAnalyzer_04_MPI_page-3.png}
\caption{The event at the end of the parton level in the $\eta$--$\phi$ plane (\code{genPartAnalyzer\_04\_MPI.pdf},
page 3; marker size grows with $\log\pt$, the \pt{} is written for $\pt>5\GeV$): 169 partons with their primordial
$k_\mathrm{T}$, coloured by the scattering they belong to. Orange and green: the ISR (1) and FSR (11) partons of
the hard process, the same as in figure~\ref{fig:04_showers}(b), with the $b$ at $91.7\GeV$ and the $\bar b$ at
$116.4\GeV$. Cyan: the 31 multiparton interactions, squares their 61 outgoing partons, triangles the 25 ISR (up)
and 56 FSR (down) emissions of their showers, soft and spread over the whole plane. Grey diamonds: the 13 beam
remnants, 11 of them at $\abseta>3.5$; 16 partons beyond $\abseta=6$ are crosses on the border. Dashed: the cones
$\Delta R<0.4$, $127.3\GeV$ around the $b$, of which $6.9\GeV$ from the underlying event (UE: MPI and remnants),
and $117.1\GeV$ around the $\bar b$ (UE $0.8\GeV$); the inset magnifies the $b$ cone, its UE partons in cyan. The
magenta star is $\vec\pt(b)+\vec\pt(\bar b)$, the recoil of the \bbbar{} pair against its radiation: $25.1\GeV$
($27.3\GeV$ before the primordial $k_\mathrm{T}$).}
\label{fig:04_display}
\end{figure}

\paragraph{Colour reconnection.}
With 32 interactions the event has 32 colour-connected systems stretched between the two remnants. Pythia8's
MPI-based colour reconnection (\code{ColourReconnection:mode} at its default 0, CP5 strength \code{range=5.176})
merges the colour lines of the softer interactions with the harder ones, which shortens the strings and reduces
the number of hadrons. It acts on colour indices, which \code{genParticles} does not store; its effect appears
only when the strings break in chapter~\ref{ch:05_Had}, and nothing in the momenta of this record changes.

\paragraph{Beam remnants, and why the record still has $13\,600\GeV$.}
Each proton gave one parton to each of the 32 interactions; what is left is the beam remnant. Pythia8 builds it so that the proton's flavour content is restored: a
valence quark taken out leaves a diquark (a $ud$ or $uu$ pair, pdg 2101, 2103 or 2203), a sea quark taken out
leaves its antiquark partner, and so on. This event has 13 remnant partons, all single quarks or antiquarks
($u$, $d$, $\bar u$, $\bar d$, $\bar s$) and no diquark: with 32 partons extracted from each proton,
including the valence $u$ at $x=0.39$ of the hard process, no valence pair was left intact. They carry
$1783.0\GeV$ of energy and $6.5\GeV$ of scalar \pt{} in total; eleven of the 13 have $\abseta>3.5$ and the
two hardest, 416 and $350\GeV$, are at $\eta=7.0$ and $-7.3$, in the beam pipe.\footnote{FWLite on
\code{step1\_GEN\_04\_MPI.root}: the status-1 partons whose mother is a beam proton (status 4). In chapter~\ref{ch:05_Had}
the same 13 partons carry status 63.} Energy and momentum are conserved exactly: the 169 final-state partons sum
to $E=13\,600.0\GeV$ and $\sum\vec\pt=0$, because the remnants take what the interactions did not, and the
primordial-$k_\mathrm{T}$ recoil is given to them. Of the $13.6\TeV$, $2862.2\GeV$ sits in the 14 partons of
the hard process and its showers, $8954.8\GeV$ in the 142 partons of the MPI and their showers, and
$1783.0\GeV$ in the remnants; the hard system has $266.8\GeV$ of scalar \pt, the MPI $207.9\GeV$, the remnants
$6.5\GeV$.

\paragraph{What reaches the two $b$ jets.}
For the jet chapters, how much MPI activity lands near the $b$ and $\bar b$? The cones of
figure~\ref{fig:04_display}, $\Delta R<0.4$ around the final $b$ ($\eta=-1.57$, $\phi=-2.08$) and $\bar b$,
answer it parton by parton. The $b$ cone holds nine partons with a vector \pt{} of $127.3\GeV$: the six of the
hard process of figure~\ref{fig:04_showers}(b) ($120.6\GeV$ of scalar \pt) and three of the \emph{underlying
event} (UE: the partons of the MPI and the remnants), $6.9\GeV$, a $\bar u$ of $3.7\GeV$ and a gluon of
$2.3\GeV$ from the initial-state showers of two MPI and a gluon of $0.9\GeV$ radiated by the $\bar c$ of a third.
That $\bar c$ itself, now $0.8\GeV$, is the one UE parton in the $\bar b$ cone, next to the $\bar b$ of
$116.4\GeV$: one soft interaction (a $\bar c$ and a gluon of $1.9\GeV$) puts \pt{} into both
jets.\footnote{\code{genPartAnalyzer\_04\_MPI.log} (the ``cone'' lines of stage 3 give the scattering of every
member), \code{ledger\_04\_MPI.txt}; the ancestry of the UE partons from FWLite (companion, section 4.6).} In the
tracker acceptance, $\abseta<2.5$, the 80 final-state partons are 13 from the hard system ($150.5\GeV$ of \pt), 66 from the MPI ($103.8\GeV$) and one remnant quark ($0.3\GeV$): the underlying
event of this collision is two fifths of the central \pt, spread over the whole $\eta$--$\phi$ plane, 5\% of
the \pt{} inside the cone of the $b$ and less than one per cent inside that of the $\bar b$.

\section{Ledger and display}

\begin{table}[htb]
\centering\small
\caption{Ledger of the stage: the record and the final state (status 1) before (chapter~\ref{ch:02_PS},
\code{ledger\_02\_PS.txt}) and after (\code{ledger\_04\_MPI.txt}) multiparton interactions and beam remnants. Regions:
tracker ($\abseta<2.5$), the gap to the HF edge, the forward calorimeter, and $\abseta>5.2$, which Geant4 is not
given (chapter~\ref{ch:09_G4}). Cones: the status-1 partons within $\Delta R<0.4$ of the final $b$ and $\bar b$
(\code{genLedger.py}), with their vector \pt.}
\label{tab:04_ledger}
\footnotesize\setlength{\tabcolsep}{3.2pt}
\begin{tabular}{lrrrrrrrrr}
\toprule
 & \multicolumn{3}{c}{chapter 02 (MPI off)} & \multicolumn{3}{c}{chapter 04 (MPI on)} & \multicolumn{3}{c}{difference} \\
\cmidrule(lr){2-4}\cmidrule(lr){5-7}\cmidrule(lr){8-10}
 & $n$ & $E$ [GeV] & \pt{} [GeV] & $n$ & $E$ [GeV] & \pt{} [GeV] & $n$ & $E$ [GeV] & \pt{} [GeV] \\
\midrule
record (all entries) & 34 & & & 745 & & & $+711$ & & \\
\quad hard process (21--29) & 4 & & & 4 & & & 0 & & \\
\quad MPI (31--39) & 0 & & & 124 & & & $+124$ & & \\
\quad ISR (41--49) & 0 & & & 190 & & & $+190$ & & \\
\quad FSR (51--59) & 15 & & & 192 & & & $+177$ & & \\
\quad primordial $k_\mathrm{T}$, remnants (61--69) & 2 & & & 64 & & & $+62$ & & \\
\midrule
final state (status 1), all partons & 11 & 13600.0 & 230.7 & 169 & 13600.0 & 481.2 & $+158$ & 0.0 & $+250.5$ \\
\quad $\abseta<2.5$ & 6 & 256.9 & 113.7 & 80 & 642.1 & 254.5 & $+74$ & $+385.2$ & $+140.8$ \\
\quad $2.5<\abseta<3.0$ & 0 & 0.0 & 0.0 & 17 & 190.5 & 22.3 & $+17$ & $+190.5$ & $+22.3$ \\
\quad $3.0<\abseta<5.2$ & 1 & 2413.9 & 114.0 & 50 & 4907.4 & 183.0 & $+49$ & $+2493.5$ & $+69.0$ \\
\quad $5.2<\abseta<5.5$ & 0 & 0.0 & 0.0 & 3 & 245.3 & 2.5 & $+3$ & $+245.3$ & $+2.5$ \\
\quad $\abseta>5.5$ & 4 & 10929.2 & 2.9 & 19 & 7614.7 & 18.9 & $+15$ & $-3314.5$ & $+16.0$ \\
\quad visible ($\abseta<5.2$) & 7 & 2670.8 & 227.8 & 147 & 5740.0 & 459.8 & $+140$ & $+3069.2$ & $+232.0$ \\
\midrule
$b$ cone $\Delta R<0.4$ & 6 & 256.9 & 113.6 & 9 & 325.3 & 127.3 & $+3$ & $+68.4$ & $+13.7$ \\
$\bar b$ cone $\Delta R<0.4$ & 1 & 2413.9 & 114.0 & 2 & 2485.9 & 117.1 & $+1$ & $+72.0$ & $+3.1$ \\
\bottomrule
\end{tabular}
\end{table}

Table~\ref{tab:04_ledger} closes the ledger of the stage against chapter~\ref{ch:02_PS}
(chapter~\ref{ch:03_MnM} only copied the file). The stage adds 158 final-state partons and $250.5\GeV$ of scalar
\pt{} but no energy: the $13\,600\GeV$ move from $\abseta>5.5$, where the two spectators of
chapter~\ref{ch:02_PS} sat with $10\,929\GeV$, into $3.0<\abseta<5.2$ ($+2493.5\GeV$, the MPI partons and their
showers at small $x$) and the tracker ($+385.2\GeV$); the visible energy doubles from 2670.8 to $5740.0\GeV$, and
$2.5<\abseta<3.0$, empty before, receives 17 partons and $190.5\GeV$. The two cone rows are the jet side of the
ledger: the $b$ cone gains $13.7\GeV$, half of it the underlying event ($6.9\GeV$) and half the other shower
draw (its hard-process partons carry $120.5\GeV$ of vector \pt{} instead of $113.6\GeV$).

Figure~\ref{fig:04_display} is the $\eta$--$\phi$ display at the end of the stage: 169 partons, each coloured
by the scattering it belongs to (the analyzer follows its mothers up to a scattering, as the FWLite count of
table~\ref{tab:04_shower}): 14 of the hard process and its showers, 142 of the MPI and their showers, 13 beam
remnants. The 14 are the partons of figure~\ref{fig:04_showers}(b) ($b$, $\bar b$, one ISR and 11 FSR partons,
same colours), moved only by the primordial $k_\mathrm{T}$: the same event, with the rest of the collision
added. The 142 are the 61 outgoing partons of the 31 interactions (a parton that radiated keeps its label) and
the 25 ISR and 56 FSR emissions of their showers.


\section{Open problems for Phase-2}



\paragraph{CP5 is a 2017 tune extrapolated to $14\TeV$ through one exponent.}
The block used by the fragment of this book, \path{Configuration/Generator/python/MCTunes2017/PythiaCP5Settings_cfi.py},
sets \code{pT0Ref=1.41}, \code{ecmPow=0.03344}, \code{coreRadius=0.7634}, \code{coreFraction=0.63} and
\code{ColourReconnection:range=5.176}; these were fitted to 7 and $13\TeV$ data (CMS-GEN-17-001; the release
records only the tune name), and the only handle that knows about the beam energy is \code{ecmPow}, which moves $\ptz$ from $1.4395\GeV$ at $13\TeV$ to
$1.4431\GeV$ at $14\TeV$. The Pythia8 default is \code{ecmPow=0.215} (Monash); CP5's value, six times smaller,
means an almost energy-independent regulator, so the MPI activity in CP5 grows with $\sqrt s$ nearly as fast as
the parton densities at small $x$ do. No $14\TeV$ measurement tests this extrapolation yet; the release contains no $14\TeV$ tune, and the directory \path{MCTunes2017/} holds only the
CP1--CP5 blocks of 2017 with their ``Up''/``Down'' variants and the two colour-reconnection alternatives
\path{PythiaCP5CR1TuneSettings_cfi.py} (\code{ColourReconnection:mode=1}, \code{BeamRemnants:remnantMode=1},
\code{pT0Ref=1.375}) and \path{PythiaCP5CR2TuneSettings_cfi.py} (\code{mode=2}, \code{pT0Ref=1.454},
\code{ecmPow=0.0555}), alternatives for systematic studies, not replacements. What could not be established:
which data entered the CP5 fit (the release records only the tune name) and whether a newer CMS tune exists
outside the release; inside it there is none, and no $14\TeV$ or Phase-2 variant of the MPI parameters.

\paragraph{The only Run-3 update is the total cross section, and most fragments do not use it.}
The release has a second directory, \path{Configuration/Generator/python/MCTunesRun3ECM13p6TeV/}, whose
\path{PythiaCP5Settings_cfi.py} differs from the 2017 file in two lines only: \code{SigmaTotal:sigmaEl}
(22.08 instead of $21.89\mb$) and \code{SigmaTotal:sigmaTot} (101.037 instead of $100.309\mb$). With
\code{SigmaTotal:mode = 0} Pythia8 takes these total and elastic cross sections as fixed numbers at whatever beam
energy is set, and derives from them the non-diffractive cross section that normalises the MPI rate, the
$\sigma_{\mathrm{ND}}=55.42\mb$ printed at $13.6\TeV$ in \code{step1\_GEN\_04\_MPI.log}. Our fragment and the
release's minimum-bias fragment \path{Configuration/Generator/python/MinBias_14TeV_pythia8_TuneCP5_cfi.py} (the
pileup library of chapter~\ref{ch:08_PU}, and the \path{RelValMinBias_14TeV} input that
\path{Configuration/PyReleaseValidation/python/relval_steps.py} mixes at \path{AVE_200_BX_25ns} for the Phase-2
workflows) both import the 2017 block, i.e.\ the $13\TeV$ cross sections; only six B-physics and soft-QCD
fragments import the $13.6\TeV$ one, and no file in the release carries $14\TeV$ values. At PU~200 a per-cent
change in the MPI activity of one minimum-bias collision is a per-cent change of the whole pileup energy in every
crossing.

\paragraph{Pileup is MPI without a hard process.}
The minimum-bias collisions of chapter~\ref{ch:08_PU} are generated with \code{SoftQCD:inelastic = on} and the
same CP5 block: their 7.0 interactions on average (\code{genstats\_08\_PU.txt}) come from the same $\ptz$,
matter profile and colour reconnection as the 31 interactions of this event, so the pileup jets of
chapter~\ref{ch:08_PU} are made of exactly the objects studied here. Table~\ref{tab:04_ledger} shows where they
go: 17 partons with $190.5\GeV$ added in $2.5<\abseta<3.0$ and 49 with $2493.5\GeV$ in $3.0<\abseta<5.2$, the
HGCal endcap and the extended Phase-2 tracker ($\abseta<4.0$), the regions Run~3 measures least precisely and
that constrained the tune the least.


